% ============================================================
% 01_introduccion.tex
% Responsable: Integrante X
% ============================================================

\section{Introducción}

% TODO: Contextualizar el problema de sincronización en sistemas multicore embebidos.
% Puntos clave a cubrir:
%   - Problema de la sincronización y exclusión mutua en sistemas embebidos multicore
%   - Limitaciones de los mecanismos software actuales (spin-locks, mutexes)
%   - Consumo energético como restricción crítica en sistemas embebidos
%   - Motivación para un enfoque hardware (C-Lock)

\subsection{Contexto}

Los sistemas embebidos multicore presentan el desafío de coordinar el acceso concurrente
a recursos compartidos. Los mecanismos de sincronización basados en software, como
\textit{spin-locks} y \textit{mutexes}, implican un costo energético significativo debido
a la espera activa (\textit{busy-waiting}).

\subsection{Motivación}

% TODO: Ampliar con datos cuantitativos del paper (overhead energético de spin-locks, etc.)

\subsection{Objetivo del Informe}

Este informe analiza el mecanismo C-Lock propuesto en \cite{clock2023}, describiendo
su arquitectura, sus ventajas sobre enfoques software, y el diseño de nuestra
simulación.
